Example \ref{example.idealisedPlacement} and Definition \ref{def.geometricPlacement}
turn the events of one Hawkes process into messages in two ways.
Under the first, Assumption \ref{assumption.idealisedBook} holds
and \eqref{eq.idealisedUpdate} is exact along the path;
under the second, the assumption holds at no $\offsetParameter$,
by Remark \ref{remark.depthBehindTheQuotes}.
We drive the two placements with the same events from the same book,
and compare the books they produce:
first through the time at which they part,
then through the distance between their mid-prices.

\begin{assumption}\label{assumption.twoPlacements}
 Let $(\arrivalTimes[]_n, \event[n], \orderSize_n)$ be the order flow of
 Definition \ref{def.orderFlowProcess},
 and let an opening book have both sides occupied
 and satisfy part \emph{iii.} of Assumption \ref{assumption.idealisedBook}.
 The \emph{idealised book} is folded from the opening book
 by the messages that Example \ref{example.idealisedPlacement} makes of the events
 $(\arrivalTimes[]_n, \event[n])$,
 and the \emph{geometric book} by the messages that
 Definition \ref{def.geometricPlacement} makes of
 $(\arrivalTimes[]_n, \event[n], \orderSize_n)$.\footnote{We
 take the best bid of either book above one tick.
 A limit buy that \eqref{eq.limitPrice} prices below one tick is placed at one tick
 by Definition \ref{def.geometricPlacement},
 which is behind the best bid unless the best bid is itself at one tick;
 the statements below that read a limit price assume the first case.}
\end{assumption}

A quantity of the idealised book carries the superscript $\idealisedBook$,
and one of the geometric book the superscript $\geometricBook$:
$\midPriceOf{\idealisedBook}_t$, $\nthBestBidSize[1,\geometricBook]_t$,
$\orderFlowContribution^{\geometricBook}$.
As in Definition \ref{def.orderFlowImbalance},
a quantity indexed by $\arrivalTimes[]_n$ is read after the $n$-th event,
and $\arrivalTimes[]_0 := 0$ indexes the opening book.
Write $\sigmaAlgebra^{\multiCountingProc}$ for the $\sigma$-algebra generated by the
arrival times and the types.
The idealised placement draws nothing beyond the events,
so the idealised book is a function of the opening book and of the path,
and it is $\sigmaAlgebra^{\multiCountingProc}$-measurable.
Let $\breachTime$ be the first arrival time at which the idealised placement would empty a
side;
it is $\sigmaAlgebra^{\multiCountingProc}$-measurable,
and the idealised book is defined on $[0, \breachTime)$.

\subsubsection{The first discrepancy}
\label{sec.firstDiscrepancy}

In this subsubsection the orders have size one, $\orderSize_n = 1$,
and the flow has no withdrawals:
$\baseIntensity_5 = \baseIntensity_6 = 0$,
and the fifth and sixth rows and columns of $\excitation$ are null.
The intensities of the two withdrawal types then vanish identically,
so that these types carry no event almost surely,
and $(\arrivalTimes[]_n, \event[n])$ is a Hawkes process on the four types of market and
limit orders.
We write $\countingProc[\limitTypes] := \countingProc[3] + \countingProc[4]$
for the count of limit orders,
$\mathbf{1}_{\limitTypes}$ for the indicator vector of the types $3$ and $4$,
and $\stationaryIntensity[\limitTypes]
:= \mathbf{1}_{\limitTypes}^{\transpose}\stationaryIntensity$.

\begin{defi}\label{def.targetOffset}
 Let the $n$-th event be a limit order, with $\arrivalTimes[]_n < \breachTime$,
 and let $\size^{\idealisedBook}_{\arrivalTimes[]_n-}$ be the size at the best quote of its
 own side in the idealised book immediately before $\arrivalTimes[]_n$.
 The \emph{target offset} of the event is
 \begin{equation}\label{eq.targetOffset}
  \targetOffset_n
  := \frac{\LOBspread^{\idealisedBook}_{\arrivalTimes[]_n-}}{\tickSizeOfLOB} - 1
  - \indicator{\size^{\idealisedBook}_{\arrivalTimes[]_n-} = \uniformDepth} .
 \end{equation}
\end{defi}

The target offset is the offset at which \eqref{eq.limitPrice} prices the order
where the idealised placement does:
at the best quote of its side when that quote holds less than $\uniformDepth$,
and one tick inside it when it holds $\uniformDepth$.
It equals $-1$ exactly when the quote holds $\uniformDepth$ and the spread is one tick:
the idealised order then trades at the opposite best quote,
and no offset reproduces it.

\begin{lemma}\label{lemma.agreementRule}
 Let the two books coincide immediately before $\arrivalTimes[]_n < \breachTime$.
 Then they coincide after the $n$-th event if and only if
 $\event[n]$ is a market type,
 or $\event[n]$ is a limit type and $\priceOffset_n = \targetOffset_n$.
 Given $\sigmaAlgebra^{\multiCountingProc}$ and $(\priceOffset_m)_{m < n}$,
 a limit event therefore leaves the books coinciding with probability
 \begin{equation}\label{eq.agreementProbability}
  \offsetParameter \, (1-\offsetParameter)^{\targetOffset_n} \,
  \indicator{\targetOffset_n \geq 0} .
 \end{equation}
\end{lemma}
\begin{proof}
 The book before the event satisfies Assumption \ref{assumption.idealisedBook},
 by Example \ref{example.idealisedPlacement}, since $\arrivalTimes[]_n < \breachTime$.
 A market buy is $(\arrivalTimes[]_n, 1, \infty, +1)$ under the geometric placement,
 and a limit buy of one share at $\bestAskPrice_{\arrivalTimes[]_n-}$ under the idealised one,
 whose market part is its whole size by Proposition \ref{prop.decompositionOfLimitOrder}.
 The best ask holds at least one share,
 so both trade one share there and rest nothing.
 Likewise for a market sell.
 For a limit order of direction $d$,
 let $\tilde{p}_n$ be the price of the idealised order.
 Then \eqref{eq.limitPrice} and \eqref{eq.targetOffset} give
 $d\,(p_n - \tilde{p}_n) = \tickSizeOfLOB\,(\targetOffset_n - \priceOffset_n)$,
 so the two messages are the same when $\priceOffset_n = \targetOffset_n$.
 Otherwise the geometric order rests one share at $p_n$, on side $d$.
 If $\targetOffset_n \geq 0$ the idealised order rests one share on side $d$ at
 $\tilde{p}_n \neq p_n$,
 and if $\targetOffset_n = -1$ it adds nothing to side $d$;
 in both cases the size resting at $p_n$ differs by one between the two books.
 The probability is \eqref{eq.priceOffset} at $\targetOffset_n$,
 the offsets being independent of the flow and of one another.
\end{proof}

\begin{defi}\label{def.firstDiscrepancy}
 The \emph{first discrepancy} $\discrepancyTime$ is the first arrival time
 $\arrivalTimes[]_n < \breachTime$
 after which the idealised and the geometric book differ in the size resting at some price,
 and $\breachTime$ if there is none.
\end{defi}

\begin{prop}\label{prop.coincidenceBeforeDiscrepancy}
 The idealised and the geometric book coincide on $[0, \discrepancyTime)$,
 and so do their mid-prices;
 there the geometric book satisfies Assumption \ref{assumption.idealisedBook}.
 Moreover, $\discrepancyTime$ is the first $\arrivalTimes[]_n < \breachTime$
 at which $\event[n]$ is a limit type and $\priceOffset_n \neq \targetOffset_n$,
 and $\breachTime$ if there is none.
\end{prop}
\begin{proof}
 By induction on $n$.
 The books coincide at the opening, neither changes between events,
 and while they coincide before an event,
 Lemma \ref{lemma.agreementRule} decides whether they coincide after it.
\end{proof}

The mid-prices may still coincide after $\discrepancyTime$:
an order resting behind the best quote of its side changes the book and not the mid-price.
The first discrepancy is therefore at most the first time at which the two mid-prices
differ.

\begin{prop}\label{prop.couplingSurvival}
 For $t \geq 0$ let
 $\cumulativeTargetOffset(t)$ be the sum of the target offsets of the limit events of
 $(0, t]$ that precede $\breachTime$,
 with $\cumulativeTargetOffset(t) := \infty$ if one of them is $-1$.
 Then for every $t \geq 0$
 \begin{equation}\label{eq.couplingSurvival}
  \Prob\big( \discrepancyTime > t \,\big\vert\, \sigmaAlgebra^{\multiCountingProc} \big)
  = \offsetParameter^{\countingProc[\limitTypes](t)} \,
  (1 - \offsetParameter)^{\cumulativeTargetOffset(t)} \,
  \indicator{\breachTime > t} ,
 \end{equation}
 with $(1-\offsetParameter)^{\infty} := 0$ and $0^{0} := 1$.
\end{prop}
\begin{proof}
 By Proposition \ref{prop.coincidenceBeforeDiscrepancy},
 the event $\{\discrepancyTime > t\}$ is the intersection of $\{\breachTime > t\}$
 with the events $\{\priceOffset_n = \targetOffset_n\}$ over the limit events of $(0,t]$.
 The first event and every $\targetOffset_n$ are
 $\sigmaAlgebra^{\multiCountingProc}$-measurable,
 the target offsets being read on the idealised book,
 while the offsets are independent of $\sigmaAlgebra^{\multiCountingProc}$ and of one
 another, by Definition \ref{def.geometricPlacement}.
 The conditional probability of the intersection is therefore the product of
 \eqref{eq.agreementProbability} over the limit events of $(0,t]$,
 which is \eqref{eq.couplingSurvival}.
\end{proof}

The law of the first discrepancy is thus a functional of the idealised path alone.
A limit order parts the books with probability $1 - \offsetParameter$
when it joins a best quote facing a spread of one tick,
with a probability that grows with the spread otherwise,
and with certainty at a full best quote facing a spread of one tick.

\begin{remark}\label{remark.couplingAtUnitOffsetParameter}
 Every factor of \eqref{eq.couplingSurvival} is at most $\offsetParameter$,
 with equality at a target offset of zero, so that
 \begin{equation*}
  \offsetParameter^{\countingProc[\limitTypes](t)} \,
  \indicator{\breachTime > t, \, \cumulativeTargetOffset(t) = 0}
  \leq \Prob\big( \discrepancyTime > t \,\big\vert\, \sigmaAlgebra^{\multiCountingProc} \big)
  \leq \offsetParameter^{\countingProc[\limitTypes](t)} .
 \end{equation*}
 At $\offsetParameter = 1$ every offset is zero almost surely,
 and $\discrepancyTime$ is $\sigmaAlgebra^{\multiCountingProc}$-measurable:
 it is the first limit event before $\breachTime$ whose target offset is not zero,
 and $\breachTime$ if there is none.
 The path does not depend on $\offsetParameter$,
 and as $\offsetParameter \uparrow 1$ the right-hand side of \eqref{eq.couplingSurvival}
 converges to $\indicator{\breachTime > t, \, \cumulativeTargetOffset(t) = 0}$,
 so that $\Prob(\discrepancyTime > t)$ converges to its value at $\offsetParameter = 1$
 by bounded convergence.
\end{remark}

\begin{exercise}\label{exercise.limitOrdersBeforeDiscrepancy}
 Show that
 $\Prob\big( \countingProc[\limitTypes](\discrepancyTime-) \geq k \big) \leq
 \offsetParameter^{k}$ for every $k \geq 0$,
 and deduce
 $\Expectation\big[ \countingProc[\limitTypes](\discrepancyTime-) \big] \leq
 \offsetParameter / (1 - \offsetParameter)$.
 Say under which condition on the idealised path the first bound is an equality.
\end{exercise}

By Remark \ref{remark.couplingAtUnitOffsetParameter},
$\Prob(\discrepancyTime > t)$ is bounded by
$\Expectation[\offsetParameter^{\countingProc[\limitTypes](t)}]$,
a functional of the counts that the exponential kernel computes in closed form.

\begin{prop}\label{prop.countGeneratingFunctional}
 Let $\multiCountingProc$ be an exponential-kernel Hawkes process,
 let $\pgfArgument : [0,\infty) \to [0,1]^{\numEventTypes}$ be piecewise constant,
 and let $\pgfExponent$ solve
 \begin{equation}\label{eq.pgfRiccati}
  \pgfExponent\derivative(r)
  = -\decay \, \pgfExponent(r)
  + \excitation^{\transpose}
  \big( \pgfArgument(r) \circ e^{\pgfExponent(r)} - \mathbf{1} \big),
  \qquad \pgfExponent(0) = 0,
 \end{equation}
 where $\circ$ is the entrywise product and the exponential acts entrywise.
 Then $\pgfExponent$ is defined on $[0, \infty)$,
 with values in the box
 $[-\branchingMatrix^{\transpose}\mathbf{1}, 0]$,
 and for $s \geq 0$ and $\forecastHorizon > 0$
 \begin{equation}\label{eq.countGeneratingFunctional}
  \Expectation \Big[
  \prod_{n \, : \, s < \arrivalTimes[]_n \leq s+\forecastHorizon}
  \pgfArgument_{\event[n]}(s + \forecastHorizon - \arrivalTimes[]_n)
  \,\Big\vert\, \sigmaAlgebra_s \Big]
  = \exp \Big( \pgfExponent(\forecastHorizon)^{\transpose} \decayedCounts(s+)
  + \int_{0}^{\forecastHorizon} \baseIntensity^{\transpose}
  \big( \pgfArgument(r) \circ e^{\pgfExponent(r)} - \mathbf{1} \big) \, dr \Big) .
 \end{equation}
 If moreover $\branchingRatio < 1$,
 the process is in its stationary version,
 and $\pgfArgument(r) = \mathbf{1}$ for $r \geq \forecastHorizon$,
 then the expectation of the product on the left is
 $\exp\big( \int_{0}^{\infty} \baseIntensity^{\transpose}
 ( \pgfArgument(r) \circ e^{\pgfExponent(r)} - \mathbf{1} ) \, dr \big)$,
 and the integral converges.
\end{prop}

The variable $r$ is the time remaining to the end of the window,
and $\pgfExponent_{\eone}$ is indexed by the exciting type:
an event of type $\eone$ begets events of type $e$ at the rate
$\excitation\subscriptee e^{-\decay s}$,
which is why the transpose appears.
At a constant $\pgfArgument$ the left-hand side of \eqref{eq.countGeneratingFunctional} is
the generating function
$\Expectation\big[ \prod_e \pgfArgument_e^{\countingProc(s+\forecastHorizon)
- \countingProc(s)} \,\big\vert\, \sigmaAlgebra_s \big]$
of the counts.

\begin{proof}
 On the face $\{\pgfExponent_e = 0\}$ of the box,
 the $e$-th entry of the right-hand side of \eqref{eq.pgfRiccati} is
 $(\excitation^{\transpose}(\pgfArgument \circ e^{\pgfExponent} - \mathbf{1}))_e \leq 0$;
 on the face $\{\pgfExponent_e = -(\branchingMatrix^{\transpose}\mathbf{1})_e\}$ it is at
 least $\decay(\branchingMatrix^{\transpose}\mathbf{1})_e
 - (\excitation^{\transpose}\mathbf{1})_e = 0$.
 The right-hand side is locally Lipschitz,
 so the solution exists on $[0,\infty)$ and stays in the box.

 Fix $s$ and $\forecastHorizon$, write $r(u) := s + \forecastHorizon - u$, and let
 \begin{equation*}
  \martingale^{\pgfArgument}_u
  := \prod_{n \, : \, s < \arrivalTimes[]_n \leq u}
  \pgfArgument_{\event[n]}\big(r(\arrivalTimes[]_n)\big)
  \cdot
  \exp \Big( \pgfExponent\big(r(u)\big)^{\transpose} \decayedCounts(u+)
  + \int_{0}^{r(u)} \baseIntensity^{\transpose}
  \big( \pgfArgument \circ e^{\pgfExponent} - \mathbf{1} \big) \Big)
 \end{equation*}
 for $s \leq u \leq s + \forecastHorizon$.
 Between events, by \eqref{eq.stateDynamics}, the logarithm of
 $\martingale^{\pgfArgument}$ has derivative
 $-\pgfExponent\derivative(r)^{\transpose}\decayedCounts
 - \decay\,\pgfExponent(r)^{\transpose}\decayedCounts
 - \baseIntensity^{\transpose}(\pgfArgument(r) \circ e^{\pgfExponent(r)} - \mathbf{1})$;
 at an event of type $e$ at $u$,
 $\martingale^{\pgfArgument}$ is multiplied by
 $\pgfArgument_e(r)\,e^{\pgfExponent_e(r)}$.
 Compensating the jumps with \eqref{eq.intensityFromState},
 $\martingale^{\pgfArgument}$ is the stochastic integral of
 $\martingale^{\pgfArgument}_{u-}(\pgfArgument_e(r) e^{\pgfExponent_e(r)} - 1)$
 against $d\countingProc(u) - \intensity(u)\,du$, summed over $e$,
 plus a drift in which the terms in $\baseIntensity$ cancel,
 \begin{equation*}
  \martingale^{\pgfArgument}_{u}
  \Big( -\pgfExponent\derivative(r) - \decay\,\pgfExponent(r)
  + \excitation^{\transpose}\big( \pgfArgument(r) \circ e^{\pgfExponent(r)} - \mathbf{1} \big)
  \Big)^{\transpose} \decayedCounts(u) \, du ,
 \end{equation*}
 which vanishes by \eqref{eq.pgfRiccati}.
 Since $\pgfExponent \leq 0$, $\decayedCounts \geq 0$ and $\pgfArgument \leq 1$,
 the local martingale $\martingale^{\pgfArgument}$ takes values in $[0,1]$,
 and it is therefore a martingale.
 Its values at $u = s + \forecastHorizon$ and at $u = s$ are the two sides of
 \eqref{eq.countGeneratingFunctional}.

 For the stationary version,
 apply \eqref{eq.countGeneratingFunctional} from $s - v$, $v > 0$,
 with the weight one on $(s - v, s]$.
 For $r \geq \forecastHorizon$, $\pgfExponent \leq 0$ and
 $e^{\pgfExponent} - \mathbf{1} \geq \pgfExponent$ give
 $\pgfExponent\derivative \geq \hurwitzMatrix^{\transpose}\pgfExponent$.
 The off-diagonal entries of $\hurwitzMatrix^{\transpose}$ are non-negative,
 so $e^{\hurwitzMatrix^{\transpose} r} \geq 0$ entrywise,
 and variation of constants gives
 $0 \geq \pgfExponent(r)
 \geq e^{\hurwitzMatrix^{\transpose}(r - \forecastHorizon)} \pgfExponent(\forecastHorizon)$,
 which decays exponentially by Proposition \ref{prop.hurwitzCriterion}.
 The law of $\decayedCounts\big((s-v)+\big)$ does not depend on $v$,
 so $\Expectation\big[\exp\big(\pgfExponent(\forecastHorizon+v)^{\transpose}
 \decayedCounts((s-v)+)\big)\big] \to 1$ as $v \to \infty$ by bounded convergence,
 while the integral decreases to a finite limit,
 its integrand being non-positive and bounded by
 $\vert \baseIntensity^{\transpose}\pgfExponent \vert$ for $r \geq \forecastHorizon$.
\end{proof}

\begin{exercise}\label{exercise.meanFromGeneratingFunction}
 Differentiate \eqref{eq.countGeneratingFunctional} at a constant
 $\pgfArgument = \mathbf{1}$ in the direction of $\unitVector_{\eone}$,
 and recover \eqref{eq.forwardMeanCounts}.
 Show first that the derivative of $\pgfExponent(\forecastHorizon)$ is
 $\meanResponseIntegral(\forecastHorizon)^{\transpose}\excitation^{\transpose}
 \unitVector_{\eone}$.
\end{exercise}

\begin{corol}\label{corol.decouplingRate}
 Let $\branchingRatio < 1$,
 and let $\pgfArgument := \mathbf{1} - (1-\offsetParameter)\mathbf{1}_{\limitTypes}$
 be constant.
 Then $\pgfExponent$ is non-increasing in each entry
 and converges to a zero $\pgfExponent_{\infty}$ of the right-hand side of
 \eqref{eq.pgfRiccati}.
 For every $s \geq 0$,
 \begin{equation}\label{eq.decouplingRate}
  \frac{1}{\forecastHorizon} \log \Expectation \big[
  \offsetParameter^{\countingProc[\limitTypes](s+\forecastHorizon)
  - \countingProc[\limitTypes](s)}
  \,\big\vert\, \sigmaAlgebra_s \big]
  \to - \decouplingRate(\offsetParameter)
 \end{equation}
 as $\forecastHorizon \to \infty$,
 where $\decouplingRate(\offsetParameter)
 := -\baseIntensity^{\transpose}
 \big( \pgfArgument \circ e^{\pgfExponent_{\infty}} - \mathbf{1} \big)$,
 and
 $\decouplingRate(\offsetParameter)
 = (1-\offsetParameter) \, \stationaryIntensity[\limitTypes]
 + O\big( (1-\offsetParameter)^{2} \big)$
 as $\offsetParameter \uparrow 1$.
\end{corol}
\begin{proof}
 Write $F$ for the right-hand side of \eqref{eq.pgfRiccati}.
 The derivative $v := \pgfExponent\derivative$ solves the linear equation
 $v\derivative = DF(\pgfExponent)\, v$,
 whose matrix
 $-\decay I + \excitation^{\transpose}\diag(\pgfArgument \circ e^{\pgfExponent})$
 has non-negative off-diagonal entries,
 from $v(0) = \excitation^{\transpose}(\pgfArgument - \mathbf{1}) \leq 0$;
 the fundamental solution of such an equation is non-negative, so $v \leq 0$.
 Being monotone and bounded, $\pgfExponent$ converges, and its limit is a zero of $F$.
 The integrand of \eqref{eq.countGeneratingFunctional} converges to
 $-\decouplingRate(\offsetParameter)$
 and $\pgfExponent(\forecastHorizon)$ stays bounded, which gives \eqref{eq.decouplingRate}.
 For the expansion,
 $e^{\pgfExponent} - \mathbf{1} \geq \pgfExponent$ gives
 $\pgfExponent\derivative \geq \hurwitzMatrix^{\transpose}\pgfExponent
 - (1-\offsetParameter)\,\excitation^{\transpose}\mathbf{1}_{\limitTypes}$,
 and the comparison of the previous proof gives
 $\pgfExponent_{\infty} \geq (1-\offsetParameter) \, \hurwitzMatrix^{-\transpose}
 \excitation^{\transpose}\mathbf{1}_{\limitTypes}$,
 so that $\pgfExponent_{\infty} = O(1-\offsetParameter)$.
 At $\offsetParameter = 1$ the zero of $F$ is $0$,
 where $D F = \hurwitzMatrix^{\transpose}$ is invertible,
 so $\pgfExponent_{\infty}$ is the branch of the implicit function theorem,
 whose derivative in $1-\offsetParameter$ at $0$ is
 $\hurwitzMatrix^{-\transpose}\excitation^{\transpose}\mathbf{1}_{\limitTypes}$.
 The derivative of $\decouplingRate$ in $1-\offsetParameter$ at $0$ is then
 $\baseIntensity^{\transpose}\big( \mathbf{1}_{\limitTypes}
 - \hurwitzMatrix^{-\transpose}\excitation^{\transpose}\mathbf{1}_{\limitTypes} \big)
 = -\decay\,\baseIntensity^{\transpose}\hurwitzMatrix^{-\transpose}\mathbf{1}_{\limitTypes}
 = \mathbf{1}_{\limitTypes}^{\transpose}(I - \branchingMatrix)^{-1}\baseIntensity$,
 which is $\stationaryIntensity[\limitTypes]$ by \eqref{eq.stationaryIntensity}.
\end{proof}

To first order in $1 - \offsetParameter$,
the bound $\Expectation[\offsetParameter^{\countingProc[\limitTypes](t)}]$
on $\Prob(\discrepancyTime > t)$ thus decays at the stationary rate of limit orders
times the probability that an offset is not zero.
The bound is attained along the paths on which every target offset is zero,
and the opening spread of one tick gives a regime in which these paths are described by the
counts alone.
Let
\begin{equation}\label{eq.touchWalks}
 \touchWalk{b}(t) := \nthBestBidSize[1]_{0} + \countingProc[3](t) - \countingProc[2](t),
 \qquad
 \touchWalk{a}(t) := \nthBestAskSize[1]_{0} + \countingProc[4](t) - \countingProc[1](t)
\end{equation}
be the \emph{touch walks},
and let $\walkExitTime$ be the first time at which either of them leaves
$[1, \uniformDepth]$.

\begin{prop}\label{prop.spreadPinnedAtOneTick}
 Let the opening spread be one tick.
 Then:
 \begin{enumerate}[label={\emph{\roman{*}}.}]
  \item on $[0, \walkExitTime)$ the idealised book keeps its opening best prices,
        its best sizes are $\touchWalk{b}$ and $\touchWalk{a}$,
        and every limit event has target offset zero;
        in particular $\walkExitTime \leq \breachTime$;
  \item on $\{\walkExitTime > t\}$,
        $\Prob(\discrepancyTime > t \mid \sigmaAlgebra^{\multiCountingProc})
        = \offsetParameter^{\countingProc[\limitTypes](t)}$, and
        \begin{equation*}
         0 \leq \Expectation\big[ \offsetParameter^{\countingProc[\limitTypes](t)} \big]
         - \Prob(\discrepancyTime > t)
         \leq \Expectation\big[ \offsetParameter^{\countingProc[\limitTypes](t)} ;
         \, \walkExitTime \leq t \big] ;
        \end{equation*}
  \item if a walk leaves $[1,\uniformDepth]$ through $\uniformDepth + 1$,
        then $\discrepancyTime \leq \walkExitTime$;
        at $\offsetParameter = 1$, $\discrepancyTime \geq \walkExitTime$,
        with equality at an exit through $\uniformDepth + 1$.
 \end{enumerate}
\end{prop}
\begin{proof}
 \emph{i.} By induction on the events before $\walkExitTime$,
 after each of which both walks lie in $[1,\uniformDepth]$.
 A market buy then finds $\touchWalk{a} \geq 2$ before it,
 takes one share and leaves the best ask in place;
 a limit sell finds $\touchWalk{a} \leq \uniformDepth - 1$ at a spread of one tick,
 joins the best ask,
 and has target offset $1 - 1 - 0 = 0$ by \eqref{eq.targetOffset}.
 Mutatis mutandis for the other two types.
 No event before $\walkExitTime$ empties a side.
 \emph{ii.} On $\{\walkExitTime > t\}$, $\cumulativeTargetOffset(t) = 0$ and
 $\breachTime > t$ by \emph{i.},
 and Proposition \ref{prop.couplingSurvival} gives the first claim;
 the bounds follow from it and from Remark \ref{remark.couplingAtUnitOffsetParameter}.
 \emph{iii.} An exit through $\uniformDepth + 1$ is a limit order at a full best quote
 facing a spread of one tick,
 whose target offset is $-1$.
 At $\offsetParameter = 1$ every offset is zero,
 so by \emph{i.} no limit event before $\walkExitTime$ parts the books.
\end{proof}

In the regime of Proposition \ref{prop.spreadPinnedAtOneTick}
the mid-price of the idealised book does not move,
and $\Expectation[\offsetParameter^{\countingProc[\limitTypes](t)}]$
is given by Proposition \ref{prop.countGeneratingFunctional}.
Once a best queue is cleared the spread opens to two ticks,
and by \eqref{eq.agreementProbability} a limit order then keeps the books coinciding
with probability $\offsetParameter(1-\offsetParameter)$ below the cap
and $\offsetParameter$ at it.

\subsubsection{The distance between the mid-prices}
\label{sec.distanceOfMidPrices}

We now let the orders of the geometric book carry the sizes of the flow,
and compare the two mid-prices.
The comparison rests on an identity that holds on any book.
For a quantity $X$ read after each event write
$\Delta_n X := X_{\arrivalTimes[]_n} - X_{\arrivalTimes[]_{n-1}}$.

\begin{defi}\label{def.outerSize}
 Let both sides of a book be occupied before and after its $n$-th event.
 The \emph{outer sizes} of the event are
 \begin{equation*}
  \outerSize^{b}_n :=
  \begin{cases}
   \nthBestBidSize[1]_{\arrivalTimes[]_{n-1}} & \text{if } \Delta_n \bestBidPrice > 0, \\
   \nthBestBidSize[1]_{\arrivalTimes[]_{n}} & \text{if } \Delta_n \bestBidPrice < 0, \\
   0 & \text{otherwise},
  \end{cases}
  \qquad
  \outerSize^{a}_n :=
  \begin{cases}
   \nthBestAskSize[1]_{\arrivalTimes[]_{n}} & \text{if } \Delta_n \bestAskPrice > 0, \\
   \nthBestAskSize[1]_{\arrivalTimes[]_{n-1}} & \text{if } \Delta_n \bestAskPrice < 0, \\
   0 & \text{otherwise}.
  \end{cases}
 \end{equation*}
 When a best price moves,
 its outer size is the size at the outer of its two values
 --- the lower of the two best bids, the higher of the two best asks ---
 read on the book in which that value is best.
\end{defi}

\begin{prop}\label{prop.touchIdentity}
 Let both sides of a book be occupied before and after its $n$-th event,
 and write $\orderFlowContribution = \orderFlowContribution^{b} + \orderFlowContribution^{a}$,
 with $\orderFlowContribution^{b}$ the first two and $\orderFlowContribution^{a}$ the last
 two terms of \eqref{eq.orderFlowContribution}.
 Then
 \begin{equation}\label{eq.touchIdentityOneEvent}
  \orderFlowContribution^{b} - \Delta_n \nthBestBidSize[1]
  = \signum(\Delta_n \bestBidPrice) \, \outerSize^{b}_n ,
  \qquad
  \orderFlowContribution^{a} + \Delta_n \nthBestAskSize[1]
  = \signum(\Delta_n \bestAskPrice) \, \outerSize^{a}_n .
 \end{equation}
 Consequently, over a window $(t-\ofiWindow, t]$ on which both sides stay occupied,
 \begin{equation}\label{eq.touchIdentity}
  \OFI_{t,\ofiWindow}
  - \big( \nthBestBidSize[1]_{t} - \nthBestBidSize[1]_{t-\ofiWindow} \big)
  + \big( \nthBestAskSize[1]_{t} - \nthBestAskSize[1]_{t-\ofiWindow} \big)
  = \sum_{n \, : \, t-\ofiWindow < \arrivalTimes[]_n \leq t}
  \big( \signum(\Delta_n \bestBidPrice) \, \outerSize^{b}_n
  + \signum(\Delta_n \bestAskPrice) \, \outerSize^{a}_n \big) .
 \end{equation}
\end{prop}
\begin{proof}
 If the best bid does not move,
 both bid indicators of \eqref{eq.orderFlowContribution} are one,
 $\orderFlowContribution^{b} = \Delta_n \nthBestBidSize[1]$,
 and both sides of the first identity vanish.
 If it rises, only the first is one:
 $\orderFlowContribution^{b} = \nthBestBidSize[1]_{\arrivalTimes[]_n}$,
 and the left-hand side is $\nthBestBidSize[1]_{\arrivalTimes[]_{n-1}} = \outerSize^{b}_n$.
 If it falls, only the second is one:
 $\orderFlowContribution^{b} = -\nthBestBidSize[1]_{\arrivalTimes[]_{n-1}}$,
 and the left-hand side is $-\nthBestBidSize[1]_{\arrivalTimes[]_n} = -\outerSize^{b}_n$.
 The ask is the same mutatis mutandis,
 and summing over the window telescopes the sizes.
\end{proof}

The identity uses Definition \ref{def.orderFlowImbalance}
and no property of the mechanism that produced the book.
When the best bid moves,
$\signum(\Delta_n \bestBidPrice)\,\outerSize^{b}_n
= \big( \tickSizeOfLOB \, \outerSize^{b}_n / \abs{\Delta_n \bestBidPrice} \big)
\, \Delta_n \bestBidPrice / \tickSizeOfLOB$,
so the right-hand side of \eqref{eq.touchIdentity} sums the moves of the two best prices in
ticks,
each weighted by its outer size per tick moved.
Under Assumption \ref{assumption.idealisedBook} that weight is $\uniformDepth$ at every
move.

\begin{corol}\label{corol.touchIdentityUnderTheAssumption}
 Under Assumption \ref{assumption.idealisedBook},
 every move of a best price is by one tick and has outer size $\uniformDepth$,
 and \eqref{eq.touchIdentity} is \eqref{eq.idealisedUpdate}
 with the residual \eqref{eq.idealisedResidual}.
\end{corol}
\begin{proof}
 That a best price moves only by one tick, onto or off a queue of $\uniformDepth$,
 follows from parts \emph{ii.} to \emph{iv.} of the assumption,
 as in Exercise \ref{exercise.idealisedUpdate}.
 The right-hand side of \eqref{eq.touchIdentity} is then
 $\uniformDepth \big( \bestBidPrice_t - \bestBidPrice_{t-\ofiWindow}
 + \bestAskPrice_t - \bestAskPrice_{t-\ofiWindow} \big) / \tickSizeOfLOB
 = 2\uniformDepth \, \Delta\midPrice_{t,\ofiWindow} / \tickSizeOfLOB$.
\end{proof}

\begin{prop}\label{prop.geometricContribution}
 Let both sides of the geometric book be occupied before and after its $n$-th event,
 a market or a limit order.
 Then
 \begin{equation}\label{eq.geometricContribution}
  \orderFlowContribution^{\geometricBook} =
  \begin{cases}
   \orderSize_n \wedge \nthBestAskSize[1,\geometricBook]_{\arrivalTimes[]_{n-1}}
   & \text{for a market buy,} \\
   -\big( \orderSize_n \wedge \nthBestBidSize[1,\geometricBook]_{\arrivalTimes[]_{n-1}} \big)
   & \text{for a market sell,} \\
   \pressure_{\event[n]} \, \orderSize_n \,
   \indicator{\priceOffset_n \leq
   \LOBspread^{\geometricBook}_{\arrivalTimes[]_n-}/\tickSizeOfLOB - 1}
   & \text{for a limit order.}
  \end{cases}
 \end{equation}
\end{prop}
\begin{proof}
 A market buy changes only the ask.
 If $\orderSize_n < \nthBestAskSize[1,\geometricBook]_{\arrivalTimes[]_{n-1}}$
 the best ask stays,
 and $\orderFlowContribution^{a} = -\Delta_n\nthBestAskSize[1,\geometricBook] = \orderSize_n$;
 otherwise it rises,
 and $\orderFlowContribution^{a} = \nthBestAskSize[1,\geometricBook]_{\arrivalTimes[]_{n-1}}$,
 however far the order walks the book.
 A limit buy rests in full on the bid,
 at or above the best bid exactly when
 $\priceOffset_n \leq \LOBspread^{\geometricBook}_{\arrivalTimes[]_n-}/\tickSizeOfLOB - 1$,
 by \eqref{eq.limitPrice}.
 Behind the best bid it changes neither best price nor the size at either,
 and $\orderFlowContribution^{\geometricBook} = 0$;
 at the best bid, $\orderFlowContribution^{b} = \Delta_n\nthBestBidSize[1,\geometricBook]
 = \orderSize_n$;
 above it, the bid rises and
 $\orderFlowContribution^{b} = \nthBestBidSize[1,\geometricBook]_{\arrivalTimes[]_n}
 = \orderSize_n$.
 Mutatis mutandis for the sell types.
\end{proof}

A limit order of the geometric book contributes its signed size when it reaches the best
quote of its side,
and nothing when it rests behind it;
Proposition \ref{prop.signedSizeContribution} is the case in which every order reaches it.

\begin{prop}\label{prop.geometricWithdrawalContribution}
 Let both sides of the geometric book be occupied before and after its $n$-th event,
 a withdrawal from the bid at the price $p_n$ drawn by
 Definition \ref{def.geometricPlacement}~\emph{iii.}
 Then
 \begin{equation}\label{eq.geometricWithdrawalContribution}
  \orderFlowContribution^{\geometricBook}
  = \pressure_{\event[n]} \,
  \big( \orderSize_n \wedge \size^{b}_{\arrivalTimes[]_n-}(p_n) \big) \,
  \indicator{p_n = \bestBidPrice_{\arrivalTimes[]_n-}} ,
 \end{equation}
 with the sizes and the best bid those of the geometric book;
 mutatis mutandis for a withdrawal from the ask.
\end{prop}
\begin{proof}
 The withdrawal removes
 $v_n := \orderSize_n \wedge \size^{b}_{\arrivalTimes[]_n-}(p_n)$ at $p_n$,
 and changes only the bid.
 Behind the best bid it changes neither best price nor the size at either,
 and $\orderFlowContribution^{\geometricBook} = 0$.
 At the best bid, either $v_n$ is less than the size there,
 the bid stays and $\orderFlowContribution^{b} = \Delta_n \nthBestBidSize[1,\geometricBook]
 = -v_n$;
 or the queue empties, the bid falls to the next occupied price,
 and $\orderFlowContribution^{b} = -\nthBestBidSize[1,\geometricBook]_{\arrivalTimes[]_{n-1}}
 = -v_n$.
 In both cases the contribution is $-v_n = \pressure_{\event[n]} v_n$.
\end{proof}

With orders of size one,
\eqref{eq.thinnedImbalance} therefore holds with withdrawals,
the sum running over the limit events placed behind the best quote of their side
and the withdrawals drawn behind it.

\begin{corol}\label{corol.thinnedImbalance}
 Let the orders have size one and the flow have no withdrawals,
 let $t < \breachTime$,
 and let no side of the geometric book be empty on $[0,t]$.
 Then
 \begin{equation}\label{eq.thinnedImbalance}
  \OFI^{\geometricBook}_{t,t} - \OFI^{\idealisedBook}_{t,t}
  = - \sum_{n \, : \, \arrivalTimes[]_n \leq t, \, \event[n] \in \{3,4\}}
  \pressure_{\event[n]} \,
  \indicator{\priceOffset_n \geq \LOBspread^{\geometricBook}_{\arrivalTimes[]_n-}/\tickSizeOfLOB} .
 \end{equation}
 At $\offsetParameter = 1$ the sum is empty,
 and both imbalances are $\pressure^{\transpose}\multiCountingProc(t)$.
\end{corol}
\begin{proof}
 Proposition \ref{prop.geometricContribution} with $\orderSize_n = 1$,
 the best quotes holding at least one share,
 and Proposition \ref{prop.signedSizeContribution} for the idealised book.
\end{proof}

\begin{prop}\label{prop.depthAccumulates}
 Let the flow have no withdrawals,
 and write $\size^{b}_t(\tickSizeOfLOB\N) := \sum_{p} \size^{b}_t(p)$
 for the size resting on the whole bid of the geometric book.
 While no side of the geometric book is empty,
 \begin{equation}\label{eq.depthAccumulates}
  \size^{b}_t(\tickSizeOfLOB\N)
  = \size^{b}_0(\tickSizeOfLOB\N)
  + \sum_{n \, : \, \arrivalTimes[]_n \leq t, \, \event[n] = 3} \orderSize_n
  - \sum_{n \, : \, \arrivalTimes[]_n \leq t, \, \event[n] = 2} \orderSize_n ,
 \end{equation}
 and if $\branchingRatio < 1$
 the increment of the right-hand side over $(s, s+\forecastHorizon]$
 has conditional mean
 $\meanSize \, (\unitVector_3 - \unitVector_2)^{\transpose}
 \big( \forecastHorizon\,\stationaryIntensity
 + \excitation\,\meanResponseIntegral(\forecastHorizon)\,(\decayedCounts(s+) - m) \big)$
 given $\sigmaAlgebra_s$.
 The ask is the same, with the types $4$ and $1$.
\end{prop}
\begin{proof}
 A limit buy rests in full on the bid,
 a market sell that leaves the bid occupied removes its size from it,
 and no other event changes the bid.
 The sizes are independent of the arrival times and of the types, with mean $\meanSize$,
 and \eqref{eq.forwardMeanCountsDeviation} gives the mean of the counts.
\end{proof}

When $\stationaryIntensity[3] > \stationaryIntensity[2]$
the geometric book therefore gains size on its bid at a constant mean rate,
and on its ask at the same rate under direction symmetry.
The idealised book holds at most $\uniformDepth$ at every price by construction,
so that the same excess of limit orders moves its best prices instead.

\begin{corol}\label{corol.geometricExhaustion}
 Let the flow have no withdrawals, or the orders have size one.
 Then:
 \begin{enumerate}[label={\emph{\roman{*}}.}]
  \item the first time at which a side of the geometric book is empty
        is a function of the events and of the sizes,
        the same at every $\offsetParameter$;
  \item if the orders have size one, $\branchingRatio < 1$ and
        $\stationaryIntensity[3] \leq \stationaryIntensity[2] + \stationaryIntensity[5]$,
        the bid of the geometric book is empty in finite time almost surely,
        and mutatis mutandis the ask.
 \end{enumerate}
\end{corol}
\begin{proof}
 \emph{i.} Without withdrawals \eqref{eq.depthAccumulates} reads the size on each side off
 the events and the sizes alone;
 with orders of size one a withdrawal from an occupied side removes one share,
 wherever it is drawn,
 and \eqref{eq.depthAccumulates} holds with the count of withdrawals from the bid
 subtracted.
 A side is empty at the first event at which the right-hand side reaches zero or below.
 \emph{ii.} With $u := \unitVector_3 - \unitVector_2 - \unitVector_5$,
 the bid is occupied at $t$ only if
 $u^{\transpose}\multiCountingProc(s) > -\size^{b}_0(\tickSizeOfLOB\N)$ for every $s \leq t$.
 By the functional central limit theorem of \cite{BDHM13som},
 $\big( \multiCountingProc(n s) - n s \, \stationaryIntensity \big)/\sqrt{n}$
 converges in law, in the Skorokhod space, to a Brownian motion of covariance
 $\countCovarianceRate$ of \eqref{eq.flowVolatility},
 restricted to the types that carry events.
 Since $u^{\transpose}\stationaryIntensity \leq 0$,
 for every $y > 0$ and $n$ large enough the probability that the bid is empty by $n$
 is at least that of
 $\inf_{s \leq 1} u^{\transpose}( \multiCountingProc(ns) - ns\,\stationaryIntensity)/\sqrt{n}
 < -y$,
 whose lower limit is at least the probability that a Brownian motion of variance
 $u^{\transpose}\countCovarianceRate u > 0$ falls below $-y$ before time one,
 the set being open.
 That probability tends to one as $y \downarrow 0$.
\end{proof}

\begin{remark}\label{remark.depthBalance}
 With sizes larger than one,
 a withdrawal takes $\orderSize_n \wedge \size^{b}_{\arrivalTimes[]_n-}(p_n)$,
 the only term of the balance of the bid that depends on the book,
 and the sizes enter that balance with the common mean $\meanSize$.
 If $\stationaryIntensity[3] > \stationaryIntensity[2] + \stationaryIntensity[5]$,
 the size on the bid is at least a walk whose mean grows linearly,
 and the depth of the geometric book has no stationary law.
 If $\stationaryIntensity[3] \leq \stationaryIntensity[2] + \stationaryIntensity[5]$
 with orders of size one,
 the bid is emptied by Corollary \ref{corol.geometricExhaustion}.
 A stationary depth therefore needs
 $\stationaryIntensity[3] < \stationaryIntensity[2] + \stationaryIntensity[5]$
 and sizes that are not all one:
 the excess of the withdrawals over the limit orders is then absorbed by the truncation
 of the withdrawals that draw a level holding less than their size,
 which acts the more often the thinner the book.
\end{remark}

\begin{prop}\label{prop.distanceDecomposition}
 Let $t < \breachTime$,
 and let no side of the geometric book be empty on $[0,t]$.
 Then
 \begin{equation}\label{eq.distanceDecomposition}
 \begin{split}
  \frac{2\uniformDepth}{\tickSizeOfLOB}
  \big( \midPriceOf{\geometricBook}_t - \midPriceOf{\idealisedBook}_t \big)
  =\;&
  \big( \OFI^{\geometricBook}_{t,t} - \OFI^{\idealisedBook}_{t,t} \big)
  \\
  &- \Big[ \big( \nthBestBidSize[1,\geometricBook]_t - \nthBestBidSize[1,\idealisedBook]_t \big)
  - \big( \nthBestAskSize[1,\geometricBook]_t - \nthBestAskSize[1,\idealisedBook]_t \big)
  \Big]
  \\
  &- \sum_{n \, : \, \arrivalTimes[]_n \leq t}
  \Big[ \signum(\Delta_n \nthBestBidPrice[1,\geometricBook])
  \Big( \outerSize^{b,\geometricBook}_n
  - \uniformDepth \, \frac{\abs{\Delta_n \nthBestBidPrice[1,\geometricBook]}}{\tickSizeOfLOB}
  \Big)
  \\
  &\hphantom{- \sum_{n \, : \, \arrivalTimes[]_n \leq t} \Big[}
  + \signum(\Delta_n \nthBestAskPrice[1,\geometricBook])
  \Big( \outerSize^{a,\geometricBook}_n
  - \uniformDepth \, \frac{\abs{\Delta_n \nthBestAskPrice[1,\geometricBook]}}{\tickSizeOfLOB}
  \Big) \Big] .
 \end{split}
 \end{equation}
\end{prop}
\begin{proof}
 Write \eqref{eq.touchIdentity} for the geometric book over $(0,t]$,
 and add and subtract
 $\uniformDepth \abs{\Delta_n \nthBestBidPrice[1,\geometricBook]} / \tickSizeOfLOB$
 and its ask counterpart inside each signed term,
 whose sum over the window is
 $2\uniformDepth \big( \midPriceOf{\geometricBook}_t - \midPriceOf{\geometricBook}_0 \big)
 / \tickSizeOfLOB$.
 Subtract Corollary \ref{corol.touchIdentityUnderTheAssumption} for the idealised book,
 and use that the two opening books coincide.
\end{proof}

The three terms of \eqref{eq.distanceDecomposition} are
the mismatch of the two imbalances,
the mismatch of the sizes at the touches,
and the mismatch of the depth through which the geometric book moves.
With its sign, the second is
$(2\uniformDepth/\tickSizeOfLOB)(\residual^{\geometricBook}_{t,t}
- \residual^{\idealisedBook}_{t,t})$,
with $\residual^{\geometricBook}_{t,t}$ the expression \eqref{eq.idealisedResidual}
evaluated on the geometric book.
The third vanishes on every move of the geometric book
that is by one tick onto or off a queue of $\uniformDepth$.
With orders of size one and no withdrawals,
the first is \eqref{eq.thinnedImbalance},
and it vanishes at $\offsetParameter = 1$.

Under direction symmetry the distance has no preferred sign.
Let $\swapMatrix$ act on the types as in Definition \ref{def.directionSymmetry},
and call the path $(\arrivalTimes[]_n, \swapMatrix(\event[n]))_n$,
with the same sizes and offsets,
its \emph{mirror}.

\begin{prop}\label{prop.mirrorSymmetry}
 Let the specification be direction-symmetric,
 let the path start from a state $z$ with $\swapMatrix z = z$,
 and let the opening book be symmetric about its mid-price $\midPrice_0$:
 $\size^{b}_0(p) = \size^{a}_0(2\midPrice_0 - p)$ for every $p$.
 Let $G_t$ be the event that,
 on the path and on its mirror,
 no side of either book is exhausted on $(0,t]$
 and no limit price meets the floor of Definition \ref{def.geometricPlacement}.
 Then $\big( \midPriceOf{\geometricBook}_t - \midPriceOf{\idealisedBook}_t \big)
 \indicator{G_t}$ has a law symmetric about zero,
 and in particular its mean is zero.
\end{prop}
\begin{proof}
 The reflection $p \mapsto 2\midPrice_0 - p$ maps the tick grid to itself,
 since $2\midPrice_0 = \bestBidPrice_0 + \bestAskPrice_0$,
 and it exchanges the bid and the ask.
 The matching rule, the messages of Example \ref{example.idealisedPlacement}
 and those of Definition \ref{def.geometricPlacement} commute with it,
 the type being replaced by its mirror:
 a limit buy at $\bestAskPrice - \tickSizeOfLOB(1+\priceOffset_n)$ reflects to a limit sell at
 the reflected best bid plus $\tickSizeOfLOB(1+\priceOffset_n)$,
 and a withdrawal drawn in proportion to the sizes of one side
 reflects to one drawn in proportion to the sizes of the other,
 so that the level of a withdrawal of the mirror may be drawn as the reflection of the
 level drawn on the path.
 By induction on $n$, on $G_t$ the books of the mirror are the reflections of the books
 of the path,
 so the distance of their mid-prices changes sign,
 and $G_t$ is invariant under the exchange of the path with its mirror.
 By \eqref{eq.directionSymmetry} the mirror has intensity
 $\swapMatrix\intensity[ ] = \baseIntensity + \excitation\swapMatrix\decayedCounts$,
 so it is a Hawkes process with the same specification, started from
 $\swapMatrix z = z$,
 and it has the law of the path;
 the marks are independent of both.
\end{proof}

\begin{prop}\label{prop.distanceSecondMoments}
 Take the moments below under the conditional law given that $t < \breachTime$
 and that no side of the geometric book is empty on $[0,t]$.
 Let
 \begin{equation*}
  \effectiveImpact_t
  := \frac{\CoVariance\big( \midPriceOf{\geometricBook}_t ,
  \pressure^{\transpose}\multiCountingProc(t) \big)}
  {\tickSizeOfLOB \, \Variance\big( \pressure^{\transpose}\multiCountingProc(t) \big)} ,
  \qquad
  \impactRemainder_t
  := \frac{\midPriceOf{\geometricBook}_t - \midPrice_0}{\tickSizeOfLOB}
  - \effectiveImpact_t \, \pressure^{\transpose}\multiCountingProc(t) ,
 \end{equation*}
 and
 $Q_t := \big( \effectiveImpact_t - 1/(2\uniformDepth) \big)^{2}
 \Variance\big( \pressure^{\transpose}\multiCountingProc(t) \big)
 + \Variance(\impactRemainder_t)$.
 Then
 \begin{equation}\label{eq.distanceSecondMoments}
  \Big\vert \, \Variance\Big(
  \frac{\midPriceOf{\geometricBook}_t - \midPriceOf{\idealisedBook}_t}{\tickSizeOfLOB}
  \Big) - Q_t \Big\vert
  \leq 1 + 2\sqrt{Q_t} .
 \end{equation}
\end{prop}
\begin{proof}
 The orders of the idealised book have size one,
 so $\OFI^{\idealisedBook}_{t,t} = \pressure^{\transpose}\multiCountingProc(t)$
 by Proposition \ref{prop.signedSizeContribution},
 so that by \eqref{eq.idealisedUpdate}
 $(\midPriceOf{\idealisedBook}_t - \midPrice_0)/\tickSizeOfLOB
 = \pressure^{\transpose}\multiCountingProc(t)/(2\uniformDepth) + r_t$
 with $\abs{r_t} \leq 1$.
 The distance in ticks is therefore $X_t - r_t$, with
 $X_t := \impactRemainder_t
 + \big( \effectiveImpact_t - 1/(2\uniformDepth) \big)\pressure^{\transpose}\multiCountingProc(t)$.
 By the choice of $\effectiveImpact_t$ the two terms of $X_t$ are uncorrelated,
 so $\Variance(X_t) = Q_t$,
 and $\Variance(X_t - r_t) - Q_t = \Variance(r_t) - 2\,\CoVariance(X_t, r_t)$,
 which the Cauchy--Schwarz inequality bounds by $1 + 2\sqrt{Q_t}$.
\end{proof}

The coefficient $\effectiveImpact_t$ is an effective impact of a unit of signed count on the
geometric mid-price, in ticks,
to be read against the impact $1/(2\uniformDepth)$ of the idealised book.
The variance of the distance is,
up to a term of the order of its square root,
the variance of the signed count weighted by the square of the difference of the two
impacts,
plus the variance of the part of the geometric mid-price that the signed count does not
explain.

\subsubsection{With withdrawals}
\label{sec.withWithdrawals}

We now restore the withdrawals.
For the coupling the orders keep size one;
for the distance the geometric book carries the sizes of the flow.
For a side of the idealised book,
$\size^{\idealisedBook}_{\arrivalTimes[]_n-}$ is the size at its best quote,
as in Definition \ref{def.targetOffset},
and $\size^{\idealisedBook}_{\arrivalTimes[]_n-}(\tickSizeOfLOB\N)$ the size resting on the
whole side.

\begin{lemma}\label{lemma.withdrawalAgreement}
 Let the orders have size one,
 let the two books coincide immediately before $\arrivalTimes[]_n < \breachTime$,
 and let the $n$-th event be a withdrawal.
 Then the books coincide after the event if and only if
 the price $p_n$ drawn by Definition \ref{def.geometricPlacement}~\emph{iii.}
 is the best price of its side.
 Given $\sigmaAlgebra^{\multiCountingProc}$ and the marks of the earlier events,
 this has probability
 \begin{equation}\label{eq.withdrawalAgreement}
  \frac{\size^{\idealisedBook}_{\arrivalTimes[]_n-}}
  {\size^{\idealisedBook}_{\arrivalTimes[]_n-}(\tickSizeOfLOB\N)} ,
 \end{equation}
 read on the side of the withdrawal,
 which is at most the inverse of the number of levels of that side.
\end{lemma}
\begin{proof}
 The side is occupied, so the geometric withdrawal is sent,
 and it cancels $1 \wedge \size_{\arrivalTimes[]_n-}(p_n) = 1$ share at $p_n$;
 the idealised one cancels one share at the best price.
 The two messages are the same when $p_n$ is the best price,
 also when the queue there holds one share and both empty it.
 Otherwise the size at the best price differs by one between the two books.
 The price is drawn in proportion to the sizes of the side,
 which before the event are those of the idealised book,
 and otherwise independently of the past.
 Every level behind the best quote holds $\uniformDepth$,
 at least the size at the best quote,
 which gives the bound.
\end{proof}

\begin{prop}\label{prop.couplingSurvivalWithWithdrawals}
 Let the orders have size one.
 Then the books coincide on $[0,\discrepancyTime)$,
 and
 \begin{equation}\label{eq.couplingSurvivalWithWithdrawals}
  \Prob\big( \discrepancyTime > t \,\big\vert\, \sigmaAlgebra^{\multiCountingProc} \big)
  = \offsetParameter^{\countingProc[\limitTypes](t)} \,
  (1 - \offsetParameter)^{\cumulativeTargetOffset(t)} \,
  \indicator{\breachTime > t}
  \prod_{n \, : \, \arrivalTimes[]_n \leq t, \, \event[n] \in \{5,6\}}
  \frac{\size^{\idealisedBook}_{\arrivalTimes[]_n-}}
  {\size^{\idealisedBook}_{\arrivalTimes[]_n-}(\tickSizeOfLOB\N)} ,
 \end{equation}
 with the conventions of Proposition \ref{prop.couplingSurvival},
 the product running over the withdrawals of $(0,t]$ that precede $\breachTime$.
\end{prop}
\begin{proof}
 By induction on $n$,
 with Lemmas \ref{lemma.agreementRule} and \ref{lemma.withdrawalAgreement}.
 Given $\sigmaAlgebra^{\multiCountingProc}$ and the marks of the events before the $n$-th,
 on the event that the books coincide before it they are the idealised book,
 which is $\sigmaAlgebra^{\multiCountingProc}$-measurable,
 and the event keeps them coinciding with the probability
 \eqref{eq.agreementProbability} or \eqref{eq.withdrawalAgreement},
 or one at a market order.
 The conditional probability that every event of $(0,t]$ does so is the product.
\end{proof}

In a book of many levels each withdrawal factor of \eqref{eq.couplingSurvivalWithWithdrawals}
is small,
and the two books part at the first withdrawal with probability close to one.
The coupling then says little,
and the comparison of the mid-prices has to be made on the two books as they are.

Over a window on which the spread of the idealised book exceeds one tick,
a limit buy contributes $+1$ to its bid,
a market sell and a withdrawal from the bid contribute $-1$,
and no other event contributes to it.
By Proposition \ref{prop.touchIdentity} and
Corollary \ref{corol.touchIdentityUnderTheAssumption},
the position $\nthBestBidSize[1,\idealisedBook] + \uniformDepth \,
\nthBestBidPrice[1,\idealisedBook] / \tickSizeOfLOB$
then moves by $(\unitVector_3 - \unitVector_2 - \unitVector_5)^{\transpose}$
times the increment of $\multiCountingProc$,
whose mean over a horizon $\forecastHorizon$ is
$\forecastHorizon \, (\stationaryIntensity[3] - \stationaryIntensity[2]
- \stationaryIntensity[5])$ up to a bounded term.
When $\stationaryIntensity[3] < \stationaryIntensity[2] + \stationaryIntensity[5]$,
the bid of the idealised book falls in mean, its ask rises, and its spread opens,
until a side is consumed at $\breachTime$.
Its mid-price does not inherit the drift:
with orders of size one,
Propositions \ref{prop.idealisedUpdate} and \ref{prop.signedSizeContribution} give
\begin{equation}\label{eq.idealisedMidAsCount}
 \frac{\midPriceOf{\idealisedBook}_t - \midPriceOf{\idealisedBook}_0}{\tickSizeOfLOB}
 = \frac{\pressure^{\transpose}\multiCountingProc(t)}{2\uniformDepth}
 + \frac{\residual^{\idealisedBook}_{t,t}}{\tickSizeOfLOB} ,
 \qquad
 \Big\vert \frac{\residual^{\idealisedBook}_{t,t}}{\tickSizeOfLOB} \Big\vert \leq 1 ,
\end{equation}
at every spread, for $t < \breachTime$.

The geometric book, on the other hand, may have a stationary depth,
by Remark \ref{remark.depthBalance}.
We read it from its best bid:
let $\bookShape_t$ be the geometric book translated so that its best bid is at zero,
and $\geometricState_t := (\decayedCounts(t), \bookShape_t)$.
For an event of type $e$ that finds the state $y$,
let $\midMove_e(y)$ be the move of the geometric mid-price in ticks,
a random variable through the marks of the event,
with $\Expectation^{\mathrm{mk}}$ the expectation over them,
and let
$\midDrift(y) := \sum_e \intensity(z) \, \Expectation^{\mathrm{mk}} \midMove_e(y)$,
with $z$ the first component of $y$ and $\intensity(z) = \baseIntensity_e + (\excitation z)_e$.
Let $\clusterPressure := (I - \branchingMatrix)^{-\transpose}\pressure$,
whose $e$-th entry is the expected signed pressure of the cluster of an event of type $e$,
so that
$\pressure^{\transpose}\countCovarianceRate\pressure
= \sum_e \stationaryIntensity[e] \, \clusterPressure_e^{2}$
by \eqref{eq.flowVolatility}.
We make three hypotheses on the geometric book:
\begin{enumerate}[label={\emph{H\arabic*}.}, start=0]
 \item \emph{translation invariance}: almost surely no side of the geometric book is ever
       empty and no limit price meets the floor,
       so that every message depends on the book through $\bookShape$ alone,
       and $\geometricState$ is a Markov process;\footnote{The reference price $p_0$ at
       which Definition \ref{def.geometricPlacement} prices a limit order when the opposite
       side is empty is fixed outside the book, and so is the floor at one tick.
       A placement that reads either is not invariant under a translation of the book,
       and $\bookShape$ no longer determines the next message.}
 \item \emph{stationarity}: $\geometricState$ is stationary;
 \item \emph{solvability}: there is a function $\poissonSolution$ with
       $\Expectation[\poissonSolution(\geometricState_0)^{2}] < \infty$ and
       $\generatorL\poissonSolution = -\midDrift$,
       $\generatorL$ the extended generator of $\geometricState$,
       and the corrected moves
       $\correctedMove_e(y) := \midMove_e(y) + \poissonSolution(y'_e) - \poissonSolution(y)$,
       with $y'_e$ the state after the event,
       satisfy
       $\Expectation\big[ \sum_e \intensity(\decayedCounts(0)) \,
       \Expectation^{\mathrm{mk}} \correctedMove_e(\geometricState_0)^{2} \big] < \infty$.
\end{enumerate}

\begin{prop}\label{prop.stationaryDistanceRate}
 Let the specification be direction-symmetric with $\branchingRatio < 1$,
 let the geometric book satisfy hypotheses H0 to H2,
 and let the idealised book open at time $0$ from a book satisfying part \emph{iii.} of
 Assumption \ref{assumption.idealisedBook}.
 Let $\effectiveImpact_t$ and $\impactRemainder_t$ be as in
 Proposition \ref{prop.distanceSecondMoments},
 with $\midPriceOf{\geometricBook}_t$ replaced by its increment
 $\midPriceOf{\geometricBook}_t - \midPriceOf{\geometricBook}_0$
 and the moments taken under $\Prob$, and let
 \begin{equation*}
  \Delta_t := \frac{\midPriceOf{\geometricBook}_t - \midPriceOf{\geometricBook}_0}{\tickSizeOfLOB}
  - \frac{\pressure^{\transpose}\multiCountingProc(t)}{2\uniformDepth} .
 \end{equation*}
 Then, as $t \to \infty$,
 \begin{equation}\label{eq.effectiveImpact}
 \begin{split}
  \effectiveImpact_t \to \effectiveImpact
  &:= \frac{\Expectation\big[ \sum_e \intensity(\decayedCounts(0)) \, \clusterPressure_e \,
  \Expectation^{\mathrm{mk}} \correctedMove_e(\geometricState_0) \big]}
  {\pressure^{\transpose}\countCovarianceRate\pressure} ,
  \\
  \frac{\Variance(\impactRemainder_t)}{t} \to \orthogonalVolatility^{2}
  &:= \Expectation\Big[ \sum_e \intensity(\decayedCounts(0)) \,
  \Expectation^{\mathrm{mk}} \big( \correctedMove_e(\geometricState_0)
  - \effectiveImpact\,\clusterPressure_e \big)^{2} \Big] ,
 \end{split}
 \end{equation}
 and
 \begin{equation}\label{eq.stationaryDistanceRate}
  \frac{\Variance(\Delta_t)}{t}
  \to \Big( \effectiveImpact - \frac{1}{2\uniformDepth} \Big)^{2}
  \pressure^{\transpose}\countCovarianceRate\pressure
  + \orthogonalVolatility^{2} .
 \end{equation}
 On $\{t < \breachTime\}$ the distance between the increments of the two mid-prices, in
 ticks, differs from $\Delta_t$ by at most one.
 Moreover
 $\orthogonalVolatility^{2} \geq \Expectation\big[ \sum_e \intensity(\decayedCounts(0)) \,
 \Variance^{\mathrm{mk}} \correctedMove_e(\geometricState_0) \big]$,
 with $\Variance^{\mathrm{mk}}$ the variance over the marks of the event.
\end{prop}
\begin{proof}
 By H2 and Dynkin's formula,
 $(\midPriceOf{\geometricBook}_t - \midPriceOf{\geometricBook}_0)/\tickSizeOfLOB
 = \martingale^{\geometricBook}_t + \poissonSolution(\geometricState_0)
 - \poissonSolution(\geometricState_t)$,
 with $\martingale^{\geometricBook}$ a square-integrable martingale
 that jumps by $\correctedMove_{\event[n]}$ at $\arrivalTimes[]_n$.
 By \eqref{eq.countIncrementIdentity} and
 $\pressure^{\transpose}\stationaryIntensity = 0$,
 $\pressure^{\transpose}\multiCountingProc(t) = \clusterPressure^{\transpose}
 (\martingale(t) - \martingale(0))
 - \pressure^{\transpose}\branchingMatrix(I - \branchingMatrix)^{-1}
 (\decayedCounts(t) - \decayedCounts(0))$,
 and the martingale term jumps by $\clusterPressure_{\event[n]}$ at $\arrivalTimes[]_n$.
 The two martingales jump at the same times,
 so the expectation of their covariation is
 $t \, \Expectation\big[ \sum_e \intensity(\decayedCounts(0)) \, \clusterPressure_e \,
 \Expectation^{\mathrm{mk}} \correctedMove_e(\geometricState_0) \big]$
 by H1,
 and the variance of $\martingale^{\geometricBook}_t$ is
 $t \, \Expectation\big[ \sum_e \intensity(\decayedCounts(0)) \,
 \Expectation^{\mathrm{mk}} \correctedMove_e(\geometricState_0)^{2} \big]$.
 The remaining terms are bounded in $L^{2}$ by H1, H2 and Theorem \ref{thm.lyapunov},
 and contribute $O(\sqrt{t})$ by the Cauchy--Schwarz inequality.
 With $\Expectation[\intensity(\decayedCounts(0))] = \stationaryIntensity[e]$
 this gives the limit of $\effectiveImpact_t$,
 and
 $\Variance(\impactRemainder_t)/t \to
 \Expectation\big[ \sum_e \intensity \, \Expectation^{\mathrm{mk}} \correctedMove_e^{2} \big]
 - \effectiveImpact^{2} \, \pressure^{\transpose}\countCovarianceRate\pressure$,
 which is $\orthogonalVolatility^{2}$ on expanding the square.
 The two terms of
 $\Delta_t = \impactRemainder_t
 + (\effectiveImpact_t - 1/(2\uniformDepth))\pressure^{\transpose}\multiCountingProc(t)$
 are uncorrelated, which gives \eqref{eq.stationaryDistanceRate},
 and \eqref{eq.idealisedMidAsCount} gives the comparison with the idealised book.
 The lower bound holds because $\effectiveImpact\,\clusterPressure_e$ does not depend on
 the marks.
\end{proof}

The rate \eqref{eq.stationaryDistanceRate} is a property of the geometric book alone:
the idealised book enters only through \eqref{eq.idealisedMidAsCount},
whatever its depth does.
It has two parts.
The first is the variance of the signed count weighted by the square of the difference
between the effective impact $\effectiveImpact$ and the impact $1/(2\uniformDepth)$ of the
idealised book;
the second is the part of the geometric mid-price that the signed count does not explain,
and the marks make it positive as soon as they change the corrected move on a set of
positive probability.
An effective impact alone governs the distance,
$\Variance(\impactRemainder_t) = o(t)$,
only if $\correctedMove_{e} = \effectiveImpact\,\clusterPressure_{e}$ almost surely.

\begin{remark}\label{remark.withdrawalsDeleted}
 Setting the baselines of the withdrawal types and their rows and columns of
 $\excitation$ to zero,
 without choosing the pair $(\branchingMatrix, \stationaryIntensity)$ anew by
 \eqref{eq.baselineFromTarget},
 may leave the limit flow above the consuming flow,
 $\stationaryIntensity[3] > \stationaryIntensity[2]$.
 Then the geometric book thickens on both sides at a constant mean rate,
 by Proposition \ref{prop.depthAccumulates}.
 At $\offsetParameter = 1$ a limit order arrives one tick from the opposite best quote,
 so that it narrows a spread of more than one tick and otherwise joins the best queue of its
 side;
 the best queues are then walks of positive drift,
 and the mid-price of the geometric book moves only when a market order clears one of them.
 Once neither is cleared again it stops,
 and the distance grows as the idealised mid-price does,
 at the rate $\pressure^{\transpose}\countCovarianceRate\pressure / (2\uniformDepth)^{2}$.
 The comparison then measures the accumulation of depth,
 and not the placement.
\end{remark}

\begin{remark}\label{remark.twoRegimes}
 The two settings make different statements provable.
 Without withdrawals the coupling is informative:
 the first discrepancy is decided by the limit orders alone,
 its conditional law is \eqref{eq.couplingSurvival},
 and at $\offsetParameter = 1$ with orders of size one the two imbalances coincide on every
 path, by Corollary \ref{corol.thinnedImbalance},
 so that the whole distance is carried by the sizes at the touches and by the depth,
 by \eqref{eq.distanceDecomposition}.
 The geometric book has no stationary depth there:
 if $\stationaryIntensity[3] \leq \stationaryIntensity[2]$ with orders of size one,
 a side is empty in finite time, by Corollary \ref{corol.geometricExhaustion},
 and if $\stationaryIntensity[3] > \stationaryIntensity[2]$ the book thickens,
 by Proposition \ref{prop.depthAccumulates},
 so that no effective impact is defined as a limit.
 With withdrawals the coupling says little,
 since each withdrawal keeps the books together with probability at most the inverse of the
 number of levels of its side, by Lemma \ref{lemma.withdrawalAgreement}.
 The geometric book may then have a stationary depth, by Remark \ref{remark.depthBalance},
 and Proposition \ref{prop.stationaryDistanceRate} gives the rate at which the distance grows
 as a property of the geometric book alone,
 under hypotheses H0 to H2;
 the depth of the idealised book drifts, and its mid-price does not inherit the drift.
\end{remark}
